% -------------------------------------------------------------
\subsection{跨环境系统集成与双阶段状态对齐闭环}
\label{subsec:integration}
% -------------------------------------------------------------

实机运行需要在同一台工作站上协同调度四通道 CAN 总线驱动、GPU 神经网络推理与多模态代码规划，并消除从显式几何控制切换至隐式模仿学习策略时的起始状态偏差。

\subsubsection{跨环境进程间二进制管道集成}

在系统软件栈中，机械臂与灵巧手底层驱动、RealSense 传感器接口及机器人学运动规划库运行于独立的硬件控制 Python 环境，而 LeRobot 与 PyTorch GPU 推理模块运行于另一套深度学习 Python 环境，两套环境的底层动态库版本彼此隔离。若在同一进程内混装两套依赖，或在每次调用策略时重新加载模型权重，均会影响周期控制稳定性并带来数秒级的启动停顿。

本系统通过常驻跨环境子进程与非阻塞二进制管道实现进程级解耦。双臂硬件控制进程独占四路独立 SocketCAN 接口（左右机械臂与左右灵巧手物理通道两两分离以避免高负载帧仲裁延迟）及 L515 相机流，不加载任何深度学习库。当需要运行动作分块策略时，硬件控制进程在后台拉起深度学习环境下的常驻推理子进程，通过标准输入输出字节流交换定长报文。每次策略会话启动前先发送单字节复位指令清空子进程内的动作队列；在随后的控制周期中，主控端直接从相机内存缓存读取最新一帧 $480\times 640\times 3$ 原始 \texttt{uint8} 彩色图像，与当前 46 维 \texttt{float32} 硬件反馈状态拼接后写入管道，子进程完成输入标准化、按需网络推理与输出反归一化后写回 46 维 \texttt{float32} 目标动作向量。管道读写全程采用基于文件描述符轮询的非阻塞分块传输并周期性检查准入令牌与人工保持标志，确保推理子进程异常时不会阻塞主控进程执行原位停机。

\subsubsection{两阶段交接处的起始状态对齐机制}

将代码生成智能体与双臂 ACT 策略串联运行时，两个模块交接处的起始状态是否落在训练分布内直接决定第二阶段成败。ACT 策略是在示范轨迹覆盖的状态邻域上学习的条件序列映射 $\pi_{\theta}(\hat{\bm{A}}_t \mid \bm{I}_t, \bm{s}_t, \bm{0})$。当智能体完成一次对折并将双层毛巾安放于垫面中央后，左右机械臂停留在靠近桌面中心的撤离位置，灵巧手处于仅张开拇指与食指的释放手型。若此时直接启动 ACT 策略，初始 46 维本体状态 $\bm{s}_0$ 与二次对折训练集首帧分布存在明显偏移，图像中双臂遮挡位置及其余三指姿态也与训练首帧不同，容易引起首个预测动作块大幅跳变。

为此，系统在智能体放布撤离与调用 \texttt{run\_act} 之间设计了三步起始状态对齐流程，完整交接顺序见图~\ref{fig:stage-handoff}。

\begin{figure}[htbp]
\centering
\begin{tikzpicture}[
  stepbox/.style={rectangle, draw=black!75, rounded corners=3pt, minimum width=2.45cm,
                  minimum height=1.15cm, align=center, font=\footnotesize},
  arrow/.style={-{Stealth[length=2.2mm]}, semithick}
]
  \node[stepbox, fill=orange!10] (s1) at (0, 0) {1. 织物状态准备\\一次对折并翻面\\移至垫面中央};
  \node[stepbox, fill=yellow!12] (s2) at (2.95, 0) {2. 双臂关节回位\\执行七轴 \texttt{move\_J}\\对齐 $\bm{q}_{L/R,\text{start}}$};
  \node[stepbox, fill=yellow!16] (s3) at (5.90, 0) {3. 五指手型对齐\\渐进复现示范首帧\\$\bm{r}_{L/R,0}^{\text{demo}}$ 并验收到位};
  \node[stepbox, fill=green!12] (s4) at (8.85, 0) {4. 双臂 ACT 接管\\独占双臂双灵巧手\\执行二次对折};
  \node[stepbox, fill=blue!10] (s5) at (11.80, 0) {5. 视觉结果复核\\交回智能体取图\\核验四层叠放状态};

  \draw[arrow] (s1) -- (s2);
  \draw[arrow] (s2) -- (s3);
  \draw[arrow] (s3) -- (s4);
  \draw[arrow] (s4) -- (s5);
\end{tikzpicture}
\caption{从代码智能体切入双臂 ACT 策略的三步起始状态对齐与视觉复核流程}
\label{fig:stage-handoff}
\end{figure}

第一步为织物位置与朝向核验。智能体释放毛巾后采集全局图像，检查双层长方形毛巾是否位于黑色工作垫中央区域，核对闭合折边与双层开口边的朝向是否符合二次对折示范规范；若边角存在局部翻卷，智能体先利用指尖拨正抚平再进入后续回位步骤。

第二步为双臂七轴起始关节回位。同一末端笛卡尔位姿在七自由度冗余机械臂上对应一维零空间解集，仅按笛卡尔位姿回位仍可能残留肘部构型差异。智能体直接读取配置文件中固化的二次对折录制起始关节向量 $\bm{q}_{L,\text{start}}, \bm{q}_{R,\text{start}} \in \mathbb{R}^7$，分别调用关节空间 \texttt{move\_J} 将左右机械臂驱动至录制初始关节角，并等待双侧位置保持标志确认到位静止。

第三步为双灵巧手示范首帧手型对齐。常规释放原语 \texttt{open\_gripper} 仅改变拇指与食指开度，无法恢复中指、无名指与小指在遥操作起始时刻的姿态；若对整个数据集各槽位分别取统计中位数拼凑目标手型，又可能破坏单次真实人手姿态中五指之间的几何协调关系。系统因此从二次对折训练集内选取双臂关节与 $\bm{q}_{\text{start}}$ 一致的示范首帧（第 36 条示范第 0 帧），提取其左右灵巧手各 16 个有效槽位实测值 $(\bm{r}_{L,0}^{\text{demo}}, \bm{r}_{R,0}^{\text{demo}})$ 封装为对齐接口 \texttt{prepare\_act\_hands}。该接口按周期逐步平滑驱动双手共 32 个主动槽位趋近目标手型，并在下发完成后主动查询最新槽位反馈，核验双侧全部有效槽位的稳态偏差落入容差范围。

完成上述三步对齐后，智能体调用 \texttt{run\_act(policy="fold")} 将双臂与双灵巧手控制权移交至 ACT 执行器。在设定的控制步数窗口内，系统锁定普通运动接口，由 ACT 执行器按 $20\,\mathrm{Hz}$ 节拍向常驻推理子进程请求 46 维目标动作，对机械臂关节目标完成任务系工作空间立方体与关节限位校验后经伺服接口下发，驱动双灵巧手完成双层毛巾捏持、抬升与二次对折。预定步数执行结束且双臂确认原位保持后，控制权交回智能体。智能体随即重新采集相机图像，依据最终画面检查毛巾是否已折成四层结构并留存于工作垫内，据此提交结束声明或报告异常状态。
